\section{系統功能說明}

\subsection{錄音功能}

使用者按下錄音按鍵後，系統進入錄音模式。Audio CODEC WM8731 會將麥克風或 Line-in 輸入的類比聲音轉換為數位音訊資料，FPGA 以 22.4 kHz 的取樣頻率接收 16 bit 音訊資料，並將資料依序寫入 SRAM 中。錄音過程中，LCD 顯示目前為 Recording 狀態，LED 則可顯示錄音進度或錄音狀態。

\subsection{SRAM 暫存功能}

錄製完成後，音訊資料會先保存在 SRAM 中。SRAM 屬於高速暫存記憶體，適合用於即時錄音資料儲存。使用者可以先播放 SRAM 中的錄音內容進行確認，若錄音結果不理想，則可以重新錄製並覆蓋原本 SRAM 中的資料。

\subsection{FLASH 儲存功能}

當使用者確認錄音內容後，系統會將 SRAM 中的音訊資料寫入 FLASH。FLASH 屬於非揮發性記憶體，即使 FPGA 關機後，資料仍可保存。因此本系統採用先暫存於 SRAM，再由使用者確認後寫入 FLASH 的方式，避免錄音失敗時直接覆蓋永久資料。

\subsection{播放功能}

使用者進入播放模式後，FPGA 會從 SRAM 或 FLASH 中讀取音訊資料，再透過 Audio CODEC WM8731 將數位音訊資料轉換為類比聲音輸出。播放過程中，LCD 會顯示目前播放來源與播放狀態，例如 Play SRAM 或 Play FLASH。

\subsection{FFT 頻譜顯示功能}

播放聲音時，FPGA 會同時擷取音訊資料進行 FFT 分析，將時域音訊訊號轉換為頻域資料。系統再將頻率能量分成 8 個頻帶，對應到 8$\times$8 LED 矩陣的 8 個欄位。每一欄的高度代表該頻帶的能量大小，因此播放音樂時，LED 矩陣會隨著不同頻率能量變化而跳動，達成類似音樂頻譜分析器的效果。